% ID: SEL-D0004-OB-E01
% ESTADO: validada. S005, 2026-11-07; tres pasadas del mismo autor y revisión matemática, didáctica y visual.
\textbf{a) Potencias y ecuación matricial.}\par\nopagebreak[4]
Antes de invertir una potencia, buscamos una simplificación calculando $A^2$:
\[
A^2=\begin{pmatrix}
1/4-5/4&-5/2+5/2\\
1/8-1/8&-5/4+1/4
\end{pmatrix}=-I_2.
\]
Entonces $A^4=(A^2)^2=(-I_2)^2=I_2$. Por otro lado,
\[
B^2=\begin{pmatrix}-1&-2\\4&-1\end{pmatrix},\qquad
B^2+I_2=\begin{pmatrix}0&-2\\4&0\end{pmatrix}.
\]
La ecuación se reduce a $I_2X=B^2+I_2$, cuya solución única es
\[
\boxed{X=\begin{pmatrix}0&-2\\4&0\end{pmatrix}.}
\]
Comprobación: al multiplicar por $A^4=I_2$ se recupera exactamente $B^2+I_2$; la dimensión de $X$ es $2\times2$.

\textbf{b) Invertibilidad de C.}\par\nopagebreak[4]
Las tres filas de $C$ suman la fila nula; por tanto son linealmente dependientes y $\det C=0$. En consecuencia, \textbf{C no tiene inversa}.

También se comprueba que $C(1,1,1)^t=(0,0,0)^t$: tiene un vector no nulo en su núcleo. Una matriz invertible no puede tenerlo, pues multiplicar por la inversa daría $(1,1,1)^t=0$.
